\documentclass[11pt,a4paper]{article}
\usepackage[a4paper,top=2.5cm,bottom=2.5cm,left=1.5cm,right=1.5cm]{geometry}
\usepackage[T1]{fontenc}
\usepackage[polish]{babel}
\usepackage[utf8]{inputenc}
\usepackage{indentfirst}
\usepackage{changepage}
\usepackage{amsmath}
\usepackage{tikz}
\usepackage{xcolor}

\title{Dokumentacja programu do analizy grafów}
\author{Olaf Mianowski \\ Julia Rombel}
\date{2026-04-21}

\begin{document}

\maketitle

\newpage

\tableofcontents

\newpage

\section{Wprowadzenie}
Zadanie polega na napisaniu dwóch aplikacji do wyznaczania grafów planarnych kolejno w C oraz w Javie. Każda z aplikacji ma wyznaczać współrzędne wierzchołków dla estetycznej wizualizacji grafu planarnego podanego w postaci listy krawędzi. Program stworzony w C obsługuje tylko powłokę tekstową użytkownika (TUI), zaś aplikacja w Javie ma pełną powłokę graficzną (GUI). Ta dokumentacja skupia się tylko na programie napisanym w języku Java.

\section{Cel programu}

Celem programu jest umożliwienie wizualizacji grafów planarnych w sposób przejrzysty i estetyczny z wykorzystaniem graficznego interfejsu użytkownika. Aplikacja pozwala na wczytanie grafu z pliku, przetworzenie go za pomocą wybranych algorytmów oraz przedstawienie wyników w formie graficznej. Istotnym założeniem jest interaktywność – użytkownik może modyfikować sposób prezentacji grafu oraz jego parametry wizualne w czasie rzeczywistym.

\section{Zakres programu}

Program umożliwia:
\begin{itemize}
    \item wczytywanie grafów z plików tekstowych oraz binarnych,
    \item wizualizację grafu planarnego w dedykowanym komponencie graficznym,
    \item interaktywną modyfikację widoku grafu (skalowanie, przesuwanie, zmiana stylu),
    \item obsługę różnych sposobów wyświetlania węzłów i krawędzi.
\end{itemize}

\section{Użycie programu}

Aplikacja napisana w języku Java wykorzystuje graficzny interfejs użytkownika oparty na bibliotece Swing. Interfejs umożliwia użytkownikowi wczytanie grafu z pliku, wybór algorytmu rozmieszczania wierzchołków oraz wizualną analizę grafu.

Główne okno programu składa się z trzech podstawowych elementów:

\begin{itemize}
    \item paska menu,
    \item panelu narzędziowego,
    \item komponentu graficznego odpowiedzialnego za rysowanie grafu.
\end{itemize}

Z poziomu paska menu użytkownik może wczytać graf zapisany w pliku tekstowym lub binarnym oraz zapisać wyniki działania programu. Wczytanie grafu powoduje utworzenie jego reprezentacji w pamięci programu.

Panel narzędziowy umożliwia wybór algorytmu wyznaczania współrzędnych wierzchołków oraz zmianę parametrów wizualizacji grafu. Po wybraniu algorytmu użytkownik może uruchomić proces obliczania współrzędnych, które następnie są wykorzystywane do narysowania grafu.

Centralną częścią interfejsu jest komponent graficzny, w którym wyświetlany jest graf. Komponent ten został zaimplementowany jako własna klasa dziedzicząca po \verb|JPanel|. Metoda \verb|paintComponent()| została nadpisana w celu ręcznego rysowania wierzchołków oraz krawędzi grafu.

Podczas pracy z programem użytkownik może:

\begin{itemize}
    \item zmieniać sposób wyświetlania wierzchołków i krawędzi,
    \item włączać lub wyłączać wyświetlanie etykiet,
    \item przesuwać widok grafu,
    \item zmieniać skalę wyświetlania.
\end{itemize}

Interfejs reaguje na zdarzenia generowane przez użytkownika, takie jak kliknięcia przycisków lub wybór opcji z menu, dzięki czemu obsługa programu jest intuicyjna i interaktywna.

\section{Architektura aplikacji Java}

Aplikacja została zaprojektowana w sposób modułowy z wykorzystaniem podejścia obiektowego. Poszczególne elementy programu zostały podzielone na klasy odpowiedzialne za różne aspekty działania systemu.

Do najważniejszych elementów architektury należą:

\begin{itemize}
    \item klasa reprezentująca graf oraz jego strukturę (np. wierzchołki i krawędzie),
    \item klasa opisująca pojedynczą krawędź grafu,
    \item klasa odpowiedzialna za implementację algorytmów wyznaczania współrzędnych,
    \item klasa odpowiedzialna za interfejs graficzny,
    \item klasa odpowiedzialna za rysowanie grafu w komponencie graficznym.
\end{itemize}

Główne okno programu zostało zaimplementowane jako klasa dziedzicząca po \verb|JFrame|. W tej klasie tworzony jest pasek menu, panel narzędziowy oraz komponent odpowiedzialny za wyświetlanie grafu.

Komponent wizualizujący graf został zaimplementowany jako klasa dziedzicząca po \verb|JPanel|. W klasie tej nadpisano metodę \verb|paintComponent()|, w której wykonywane jest rysowanie wierzchołków oraz krawędzi grafu przy użyciu obiektu \verb|Graphics|.

Algorytmy rozmieszczania wierzchołków zostały zaimplementowane w oddzielnych klasach, co umożliwia łatwe dodawanie nowych metod wizualizacji grafów w przyszłości. Panel narzędziowy umożliwia wybór jednego z dostępnych algorytmów oraz uruchomienie procesu obliczeń.

Taki podział programu zwiększa czytelność kodu oraz ułatwia jego wykorzystanie.

\section{Obsługa błędów}

Aplikacja zostanie wyposażona w mechanizmy obsługi błędów, które mają na celu zwiększenie stabilności programu oraz poprawę komfortu użytkownika.

Podczas wczytywania grafu program sprawdza poprawność danych wejściowych. W przypadku wykrycia nieprawidłowego formatu pliku lub niespójności danych użytkownik zostaje poinformowany o problemie za pomocą komunikatu dialogowego.

Obsługa błędów została zaimplementowana z wykorzystaniem mechanizmu wyjątków dostępnego w języku Java. W szczególności zastosowano konstrukcję \verb|try-catch|, która pozwala na przechwycenie potencjalnych błędów podczas odczytu plików, przetwarzania danych oraz działania algorytmów.

W przypadku wystąpienia błędu program wyświetla komunikat informacyjny, który pozwala użytkownikowi zidentyfikować przyczynę problemu. Dzięki temu możliwe jest szybkie podjęcie odpowiednich działań, takich jak ponowne wczytanie pliku lub zmiana parametrów programu.

Dodatkowo zastosowano zabezpieczenia zapobiegające wykonaniu operacji, które mogłyby doprowadzić do nieprawidłowego działania aplikacji, na przykład uruchomienia algorytmu bez wcześniej wczytanego grafu.

\section{Implementacja programu}

Program został zaimplementowany jako aplikacja okienkowa służąca do wizualizacji oraz analizy grafów planarnych. Aplikacja została napisana w języku Java z wykorzystaniem biblioteki Swing odpowiedzialnej za obsługę graficznego interfejsu użytkownika. Program umożliwia wczytywanie grafów z plików tekstowych oraz binarnych oraz interaktywną wizualizację grafu.

Architektura programu została podzielona na moduły odpowiedzialne za reprezentację grafu, implementację algorytmów rozmieszczania wierzchołków, obsługę interfejsu graficznego oraz operacje wejścia i wyjścia. Zastosowany podział zwiększa czytelność kodu oraz ułatwia dalszy rozwój aplikacji.

\subsection{Wczytywanie danych}

Program umożliwia wczytywanie grafów z plików tekstowych oraz binarnych. Odczyt danych realizowany jest przez klasy odpowiedzialne za obsługę wejścia i wyjścia plikowego. Podczas wczytywania pliku tekstowego każda linia interpretowana jest jako pojedyncza krawędź grafu, na podstawie której tworzona jest reprezentacja grafu w pamięci programu. Program sprawdza poprawność danych wejściowych oraz zabezpiecza się przed próbą wczytania niepoprawnego formatu pliku.

\subsection{Reprezentacja grafu}

Graf przechowywany jest w pamięci programu w postaci struktur obiektowych reprezentujących wierzchołki oraz krawędzie. Każdy wierzchołek posiada własny identyfikator oraz współrzędne wykorzystywane podczas procesu wizualizacji. Do przechowywania danych wykorzystano dynamiczne struktury danych umożliwiające efektywny dostęp do elementów grafu oraz wygodne wykonywanie akcji na jego strukturze. Pozwala to na sprawne zarządzanie wierzchołkami i krawędziami.

Podstawowe klasy odpowiedzialne za reprezentację grafu obejmują:

\begin{itemize}
    \item klasę reprezentującą cały graf,
    \item klasę reprezentującą pojedynczy wierzchołek,
    \item klasę reprezentującą pojedynczą krawędź.
\end{itemize}

\subsection{Wizualizacja grafu}

Wyświetlając graf, program pokazuje wierzchołki jako okręgi, a krawędzie są przedstawiane jako odcinki pomiędzy wierzchołkami. Użytkownik może w ustawieniach zmieniać poszczególne cechy graficznej reprezentacji grafu, takie jak:
\begin{itemize}
	\item kolor krawędzi (za pomocą menu \verb|JColorChooser|),
	\item grubość krawędzi (za pomocą suwaka \verb|JSlider|),
	\item widoczności etykiet na wierzchołkach i krawędziach (mając 3 opcje: tak, tylko po najechaniu myszką oraz nie. Opcje wybieramy za pomocą komponentu \verb|JRadioButton|).
\end{itemize}

\subsection{Obsługa zdarzeń}

Program musi reagować na wiele zdarzeń, zwłaszcza na te dotyczące opsługi plików oraz zmiany ustawień. Zdarzenia są spowodowane interakcją użytkownika z programem. Zaliczamy do nich:
\begin{itemize}
	\item rozszerzenie okna aplikacji,
	\item wybór pliku wejściowego lub wyjściowego,
	\item zamknięcie okna aplikacji,
	\item wprowadzenie zmian w wizualizacji grafu,
	\item dodanie, usunięcie lub modyfikacja wierzchołka albo krawędzi.
\end{itemize}

\section{Podsumowanie}

Projekt stanowi kompletną aplikację do wizualizacji grafów planarnych z wykorzystaniem technologii Java SWING. System został podzielony na moduły odpowiadające za logikę, interfejs użytkownika oraz przetwarzanie danych. Zastosowano i zaimplementowano podejście obiektowe, wzorce projektowe, mechanizmy obsługi błędów oraz spójności interfejsu, które zwiększają niezawodność i komfort użytkowania. Aplikacja zapewnia intuicyjną obsługę dzięki graficznemu interfejsowi użytkownika oraz umożliwia interaktywną pracę z grafem, w tym jego edycję i wizualną analizę. Program został zaprojektowany z myślą o łatwym użyciu, co pozwala większości osobom użyć go bez problemów.

\end{document}
